% !TEX root = ../main.tex
\documentclass[../main.tex]{subfiles}
\graphicspath{{\subfix{../images/}}}

\begin{document}
	
\subsection{Цели и подходы к тестированию}

В процессе разработки проводилось функциональное тестирование игровых механик, а также тестирование производительности и стабильности при работе с большим количеством снарядов. Основные цели:
\begin{itemize}
	\item проверка корректности движения игрока и обработки ввода;
	\item проверка коллизий и системы неуязвимости;
	\item проверка работы планировщика событий и правильности последовательности атак;
	\item проверка удаления снарядов за границами мира;
	\item оценка производительности (частоты кадров) при максимальной нагрузке.
\end{itemize}

Для изолированной проверки отдельных механик были созданы специальные тестовые уровни, которые становятся видимыми в меню при включённом режиме отладки (клавиша \texttt{F3}).

\subsection{Тестовые уровни}

\subsubsection{Уровень TEST}

Уровень \texttt{TEST} предназначен для проверки базовых возможностей системы снарядов. Он генерирует три типа снарядов:
\begin{itemize}
	\item две движущиеся пули (разного радиуса), вылетающие из точки (700,700) вверх;
	\item одну статическую пулю большого радиуса (для проверки коллизий в статике).
\end{itemize}
На этом уровне игроку даётся бессмертие (\texttt{player.immortal = true}), что позволяет безопасно проверять корректность отрисовки и коллизий без риска проигрыша.

\subsubsection{Уровень TEST\_HP}

Уровень \texttt{TEST\_HP} предназначен для проверки системы здоровья и неуязвимости. На нём с интервалом 1 секунду создаётся одна движущаяся пуля. Игрок имеет 3 единицы здоровья и не имеет бессмертия. Этот уровень позволяет убедиться, что:
\begin{itemize}
	\item при попадании здоровье уменьшается на 1;
	% \item после попадания игрок становится неуязвимым на 0,5 секунды (повторные попадания в этот интервал не наносят урон);
	\item при достижении здоровья 0 игра переходит в состояние \texttt{GAME\_OVER}.
\end{itemize}

\subsubsection{Уровень EMPTY}

Пустой уровень, на котором не создаётся ни одного снаряда. Используется для проверки работы игрового цикла в отсутствие атак, а также для тестирования интерфейса выбора уровней и перехода между состояниями.

\subsection{Тестирование планировщика событий}

Для проверки корректности работы событийной модели на уровнях «Beginning» и «Way» использовались следующие методы:
\begin{itemize}
	\item визуальный контроль последовательности атак (соответствие задуманным паттернам);
	\item проверка времени появления каждого события с использованием отображения \texttt{level\_time} в режиме отладки;
	\item проверка динамического добавления событий (например, при взрыве бомбы на уровне «Way»);
	\item проверка завершения уровня: после обработки всех событий и исчезновения всех снарядов игра должна перейти в состояние \texttt{LEVEL\_COMPLETED}.
\end{itemize}

В процессе тестирования были выявлены и исправлены ошибки, описанные в таблице~\ref{tab:difficulties} главы 3. В частности, проблема сортировки событий после динамического добавления была решена введением функции \texttt{sort\_attack\_events}, вызываемой после каждого добавления новых событий.

\subsection{Тестирование производительности}

Одним из ключевых требований к игре в жанре bullet hell является поддержание стабильной частоты кадров (не менее 60 FPS) при большом количестве одновременно активных снарядов. В процессе тестирования использовался встроенный счётчик FPS, доступный в режиме отладки. Кроме того, для оценки нагрузки на систему был добавлен счётчик активных снарядов.

Результаты тестирования на уровне <<Way>> показали, что на оборудовании с процессором Intel Core i5 и интегрированной графикой частота кадров оставалась не ниже 60 FPS при количестве снарядов до 300.
% основное узкое место — обход всех снарядов в цикле для проверки коллизий и удаления (линейная сложность O(n)). Дальнейшая оптимизация возможна за счёт использования пространственного разбиения (квадродерево или сетка), но для текущего объёма снарядов это не требуется.

\subsection{Режим отладки как инструмент тестирования}

Режим отладки (включается клавишей \texttt{F3}) предоставляет разработчику следующие возможности:
\begin{itemize}
	\item отображение FPS и количества активных снарядов;
	\item отображение точных координат игрока и текущего времени уровня;
	\item визуализация хитбоксов: для игрока рисуется прямоугольник;
	\item отображение счётчика попаданий (сколько раз игрок столкнулся со снарядами за уровень);
	\item доступ к тестовым уровням в меню.
\end{itemize}

Использование режима отладки позволяет быстро выявлять ошибки коллизий, проверять точное время срабатывания событий и анализировать производительность.

\subsection{Результаты тестирования}

По итогам тестирования можно сделать следующие выводы:
\begin{enumerate}
	\item Все базовые механики (движение, коллизии, неуязвимость, здоровье) работают корректно.
	\item Событийная система обеспечивает точное воспроизведение запланированных атак на уровнях «Beginning» и «Way».
	\item Производительность достаточна для комфортной игры на современном оборудовании.
	\item Режим отладки и тестовые уровни упрощают процесс выявления и исправления ошибок.
	\item Выявленные в процессе разработки трудности успешно преодолены.
\end{enumerate}

\end{document}
